\documentclass[11pt]{article}
\usepackage[margin=1in]{geometry}
\usepackage{amsmath,amssymb,amsthm}
\usepackage{graphicx}
\usepackage[hidelinks]{hyperref}

\newtheorem{theorem}{Theorem}
\newtheorem{remark}{Remark}

\newcommand{\B}{\mathbb B}
\newcommand{\C}{\mathbb C}
\newcommand{\D}{\mathbb D}

\title{Coordinate-Axis Points of the Complex Unit Ball\\Are Not MDM Points}
\author{Run \texttt{fa\_banach\_001}, lane 8}
\date{August 11, 2026}

\begin{document}
\maketitle

\section*{Status and source problem}

\noindent\textbf{Status:} candidate counterexample, likely valid; human review
requested.

In Definition~7 of their paper, PDF page~3, Kazaniecki and Wojciechowski
introduce local and global MDM points and then conjecture that the
polynomial-hull obstruction described there is the only obstruction to the MDM
property \cite{KW}.  They add that even for the Euclidean unit ball of
\(\C^2\), they do not know whether any MDM point exists.  Here is the source
passage.

\begin{center}
  \includegraphics[width=\textwidth,height=.56\textheight,keepaspectratio]
  {figures/open_problem_crop.png}
\end{center}

For clarity, put
\[
 \B^2=\{(z_1,z_2)\in\C^2:|z_1|^2+|z_2|^2<1\}.
\]
In the source definition, a boundary point \(p\) is \emph{local MDM} if some
Euclidean ball \(B(p,R)\) has the following property: for every \(c>0\), an
analytic polynomial \(F\) has all five norms
\[
 \|F\|,\quad\|F_1\|,\quad\|F_2\|,
 \quad\|F_{11}\|,\quad\|F_{22}\|
\]
at most one on \(\B^2\cap B(p,R)\), while \(|F_{12}(p)|\ge c\).
The global definition replaces the cap by all of \(\B^2\).  Subscripts denote
complex coordinate derivatives.

\section*{Idea of the proof}

At \(p=(1,0)\), restrict \(F_2\) to the axial disc and write
\(h(\zeta)=F_2(\zeta,0)\).  The transverse fiber of \(\B^2\) above a real
point \(r\) is the disc of radius \(\sqrt{1-r^2}\).  Cauchy's estimate on that
fiber, applied to \(F_{11}\), yields
\[
 |h''(r)|\le \frac1{\sqrt{1-r^2}}.
\]
This singularity is integrable at \(r=1\), so \(h'(1)=F_{12}(1,0)\) is
bounded.  Near \(p\), every sufficiently late full transverse fiber lies in
an arbitrary prescribed cap, which upgrades the global argument to a local
one.  Phase rotations and coordinate interchange then give all coordinate-axis
boundary points.

\section*{Counterexample theorem}

\begin{theorem}\label{thm:main}
For every \(\theta\in\mathbb R\), neither
\[
 p_\theta=(e^{i\theta},0)
 \qquad\text{nor}\qquad
 q_\theta=(0,e^{i\theta})
\]
is a local or global MDM point of \(\B^2\).

More quantitatively, if \(F\) is a polynomial satisfying
\[
 \|F_2\|_{L^\infty(\B^2)}\le1,
 \qquad
 \|F_{11}\|_{L^\infty(\B^2)}\le1,
\]
then
\[
 |F_{12}(1,0)|\le 1+\frac\pi2.
\]
For every \(R>0\), there is also a finite constant \(C_R\) such that the same
endpoint derivative is at most \(C_R\) whenever these two bounds hold on
\(\B^2\cap B((1,0),R)\).
\end{theorem}

\begin{proof}
We first prove the global estimate at \(p=(1,0)\).  Define
\[
 h(\zeta)=F_2(\zeta,0),\qquad \zeta\in\D.
\]
Then \(|h|\le1\) on \(\D\), and Cauchy's derivative estimate on discs of
radius tending to one gives
\[
 |h'(0)|\le1. \tag{1}
\]
For each real \(0\le r<1\), the function
\[
 g_r(w)=F_{11}(r,w)
\]
is holomorphic and bounded by one on
\(|w|<\sqrt{1-r^2}\).  Since
\(g_r'(0)=F_{112}(r,0)=h''(r)\), another Cauchy estimate gives
\[
 |h''(r)|\le \frac1{\sqrt{1-r^2}}. \tag{2}
\]
Integrating along the real radius and using (1)--(2),
\[
 |h'(r)|
 \le 1+\int_0^r\frac{dt}{\sqrt{1-t^2}}
 \le 1+\frac\pi2.
\]
Because \(F\) is a polynomial, \(h'\) is continuous at one.  Letting
\(r\uparrow1\) proves
\[
 |F_{12}(1,0)|=|h'(1)|\le1+\frac\pi2. \tag{3}
\]
The source normalization is stronger than the two hypotheses used in (3), so
\((1,0)\) cannot be a global MDM point.

We next establish the cap-local assertion.  Fix \(R>0\), and suppose only
that \(|F_2|\le1\) and \(|F_{11}|\le1\) on
\[
 U_R=\B^2\cap B((1,0),R).
\]
Set
\[
 \varepsilon=\min\left\{\frac18,\frac R4,\frac{R^2}{8}\right\},
 \qquad r_0=1-\varepsilon,
 \qquad \rho=\frac\varepsilon2.
\]
If \(|\zeta-r_0|<\rho\), then
\[
 |\zeta|<1,
 \qquad
 |1-\zeta|<\frac{3\varepsilon}{2}<R.
\]
Thus this entire axial disc lies in \(U_R\).  Cauchy's estimate for \(h\) at
its center gives
\[
 |h'(r_0)|\le\frac1\rho=\frac2\varepsilon. \tag{4}
\]

For \(r_0\le r<1\) and \(|w|<\sqrt{1-r^2}\), one has
\[
 |(r,w)-(1,0)|^2
 =(1-r)^2+|w|^2
 <(1-r)^2+(1-r^2)
 =2(1-r)
 \le2\varepsilon<R^2.
\]
Consequently the \emph{full} transverse fiber above every such \(r\) lies in
\(U_R\), and the proof of (2) applies unchanged.  From (4),
\begin{equation}
 \begin{aligned}
 |F_{12}(1,0)|
  &=|h'(1)|\\
  &\le \frac2\varepsilon+
    \int_{r_0}^1\frac{dt}{\sqrt{1-t^2}}
   =\frac2\varepsilon+\arccos(r_0)=:C_R<\infty.
 \end{aligned}
 \tag{5}
\end{equation}
Since (5) holds for every \(R>0\), no ball centered at \((1,0)\) can support
the arbitrarily large mixed derivatives required by the local MDM definition.

Finally, if \(p_\theta=(e^{i\theta},0)\), apply the result at \((1,0)\) to
\[
 G(z_1,z_2)=F(e^{i\theta}z_1,z_2).
\]
The ball and corresponding caps are preserved, all relevant norms are
unchanged, and \(|G_{12}(1,0)|=|F_{12}(p_\theta)|\).  This proves the claim for
\(p_\theta\).  Interchanging the two coordinates converts the hypotheses used
above into the source bounds on \(F_1\) and \(F_{22}\), and proves the claim
for every \(q_\theta\).
\end{proof}

\section*{Why this disproves the proposed obstruction criterion}

The closed Euclidean unit ball is its own polynomial hull.  Indeed, if
\(z\notin\overline{\B^2}\), then the holomorphic linear polynomial
\[
 L_z(w)=\overline{z_1}w_1+\overline{z_2}w_2
\]
satisfies
\[
 |L_z(z)|=\|z\|_2^2>
 \|z\|_2=\sup_{w\in\overline{\B^2}}|L_z(w)|.
\]
Thus \(p_\theta,q_\theta\) lie on, not in the interior of, the polynomial
hull.  Theorem~\ref{thm:main} nevertheless says that they are not MDM points.
This is a direct counterexample to the conjecture that hull interiority is the
only obstruction.

\section*{Verification notes}

The proof is analytic and does not depend on computation.  The reusable script
\path{code/verify_axis_bound.py} checks the antiderivative in (3), the
choice of \(\varepsilon\), and all strict cap inequalities for 15 positive
radii ranging from \(10^{-6}\) to \(10^3\), together with several intermediate
values.  It reports
\[
 1+\pi/2=2.570796326795\ldots
\]
and that all transcription checks pass.  The command used was
\begin{verbatim}
conda run --no-capture-output -n sandbox python code/verify_axis_bound.py
\end{verbatim}
These checks are not a substitute for the Cauchy-estimate proof.

\section*{Limitations and failed full-classification upgrade}

The theorem completely classifies the coordinate-axis boundary orbit under
phase rotations and coordinate interchange, but it does not decide whether a
generic boundary point \((a,b)\), with \(ab\ne0\), is MDM.  Along the analogous
coordinate slice at such a point, the available transverse radius is only of
order \(1-r\).  Direct Cauchy estimation therefore produces a
\((1-r)^{-1}\) singularity, which is not integrable.  A general unitary rotation
does not repair this because the source normalization controls the
coordinate-pure derivatives \(F_{11},F_{22}\), a family not preserved by
general unitary mixing.  Accordingly, this packet does not claim that the
unit ball has no MDM points.

\section*{Bounded novelty check}

The search performed through August 11, 2026 covered the run's registry,
solution, attempt, and proof-gap indexes; the local source corpus; and bounded
official-arXiv searches using \nolinkurl{arXiv:2301.09526}, the exact title, both author
names, ``MDM point'', ``Mitiagin-DeLeeuw-Mirkhil'', and ``unit ball''.  No
duplicate packet or later explicit resolution of this unit-ball question was
found.  This is a bounded novelty check, not a claim of priority.

\section*{Human review recommendation}

A reviewer should check (i) the identification
\(h''(r)=F_{112}(r,0)\), (ii) the two Cauchy estimates, (iii) the full-fiber cap
inclusion preceding (5), and (iv) that the source's phrase ``the above obstacle
is the only one'' is intended to cover the global MDM notion.  The
generic-boundary limitation above should remain explicit.

\begin{thebibliography}{9}

\bibitem{KW}
K. Kazaniecki and M. Wojciechowski,
\emph{On the analytic version of the Mitiagin-DeLeeuw-Mirkhil
non-inequality on bi-disc}, arXiv:2301.09526, 2023.

\end{thebibliography}

\end{document}
